\chapter{A new subject}
\label{ch:subject}

\section{The gap between two mature fields}

Reconfigurable computing classically means: a designer writes RTL, a synthesizer
emits LUTs and flip-flops, a bitstream configures an FPGA, and the fabric
evaluates the circuit at a clock~\cite{harris2012digital,gajski2009embedded}.
The fabric is spatially wired. The configuration is a bitstream. The values
on the wires are bits.

Fully homomorphic encryption classically means: a scheme evaluates circuits
\emph{on ciphertexts}~\cite{chillotti2020tfhe,ducas2015fhew}.
The ``fabric'' is an arithmetic library---NTT butterflies, key-switching
decompositions, bootstrapping. The configuration is a parameter set. The
values on the wires are polynomials in $\mathbb{Z}_q[X]/(X^N+1)$.

Each field has textbooks.
The composition does not.

\begin{definition}[Reconfigurable homomorphic computing]
The study of compiling \emph{already-synthesized} sequential netlists onto
\emph{tensor engines} so that LUT evaluation, flip-flop clocking, and
(when claimed) torus FHE occupy the same graph object, with explicit
certificates of what the graph does and does not prove.
\end{definition}

HELUT is one laboratory realization of that definition: Yosys JSON in,
\texttt{MPSGraph} out, host posedge clock, LWE/GLWE samples on the FHE
path~\cite{helut-paper}.
The subject is larger than the prototype.
A future chapter may describe a CUDA fabric, an MLIR fabric, or an ASIC LUT
plane. The definition does not name Apple.

\section{Three questions that open the course}

\begin{enumerate}
  \item \textbf{Representation.}
        Can a commodity ML graph API host exact modular arithmetic at
        $q=2^{32}$, or does the API silently round?
  \item \textbf{Clocking.}
        Can sequential logic---enables, sync resets, a RISC-V core---tick
        on that graph without the host becoming a bit-banged FPGA?
  \item \textbf{Honesty.}
        When the same object grows a real bootstrapping key, what
        certificates must ship with every tick, and which sentences remain
        forbidden?
\end{enumerate}

Pillar~I is the constructive answer to (1)--(3).
Pillar~II asks a fourth question the FPGA curriculum never asked:
\emph{what if the INIT table is not binary?}
Pillar~III asks a fifth: \emph{what if the circuit under evaluation is allowed
to mutate because the attacker in the next room is the same compiler?}

\section{Why ``reconfigurable'' is the right word}

An FPGA is reconfigurable because the LUT INIT bits and the interconnect are
late-bound.
HELUT is reconfigurable in the same sense, with three substitutions:

\begin{center}
\begin{tabular}{@{}lll@{}}
\toprule
Classical FPGA & This course & Consequence \\
\midrule
LUT INIT bitstream & Yosys \texttt{\$lut} truth table & Same IR \\
Routing fabric & Tensor edges in one graph & Place-and-route becomes lowering \\
Clock tree & Host \texttt{graph.run} loop & Timing closure becomes tick protocol \\
IOB / BRAM & Placeholders and MTLBuffers & Memory is unified, not dual-port \\
Encrypted? & Optional; certificate-gated & Oracle $\neq$ FHE \\
\bottomrule
\end{tabular}
\end{center}

The last row is the pedagogical load-bearing wall.
Students who have built FPGAs want to say ``the CPU is encrypted'' the moment
PicoRV32 ticks.
The boolean oracle path \emph{does} clock PicoRV32 on ciphertext-\emph{shaped}
vectors (\cid{1}).
The FHE path \emph{does} fetch NOPs on encrypted PicoRV32 at demo $N{=}8$
(\cid{47} CPU, \cid{51} Metal): the PC walks, decrypted $Q$ matches cleartext,
and calibrated hardness is $4.0$ bits---not the production $175.7$ row.
Covering noisy BK at $N{=}1024$ on that core is not claimed.
Chapter~\ref{ch:certs} exists so that confusion is a grading event rather than
a press cycle.

\section{The three pillars as a curriculum, not a brand}

Roadmap language names them the Turing pillar, the Schneier pillar, and a
grand challenge on proving circuits~\cite{helut-trajectory}.
In this course they are simply I, II, III, in the order a student can
verify them:

\begin{itemize}
  \item Without I, II is a float32 netlist toy and III has nowhere to run.
  \item Without II, I is a compiler course with FHE homework.
  \item Without III, I and II only attack the past; they do not design the
        next cipher under fire.
\end{itemize}

Edition~\corpusedition{} teaches I in lecture depth, Pillar~II
Theorem~\ref{thm:tensorlut} (\cid{19}) in lecture depth with empirics in
seminar, Pillar~III in seminar depth, and the grand-challenge / side-channel
items as labeled frontier (Chapter~\ref{ch:frontier}).

\section{Related systems, scoped not ranked}

\begin{center}
\begin{tabular}{@{}lcccc@{}}
\toprule
System & Exact $q{=}2^{32}$ Metal & Netlist clock & BK \texttt{\$lut} FHE & Diff.\ hardware \\
\midrule
Concrete / tfhe-rs & -- & compiler IR & yes & -- \\
OpenFHE / HElib & -- & circuit API & yes & -- \\
HELUT boolean oracle & yes & yes & -- & -- \\
HELUT encrypted path & yes & yes & yes & -- \\
TensorLUT (Pillar~II) & -- & Yosys INIT & -- & yes \\
\bottomrule
\end{tabular}
\end{center}

This is a \emph{scope} matrix, not a throughput bake-off~\cite{helut-paper}.
The encrypted-path ``yes'' is GGSW \texttt{\$lut} FHE as certified
(\cid{4}--\cid{6}, covering \cid{52}--\cid{54}); it is not PicoRV covering
at production $N$.
Concrete is a production TFHE compiler; HELUT is a netlist-clocked tensor
fabric with an FHE path.
Students comparing ``ms/gate'' across those rows without naming encoding,
$N$, and certificate status fail the five-cell test.

\begin{exercise}
Write four sentences, one per row of the matrix, that a journalist could
print without creating a \nid{}.
Then write the one-sentence lie each row most tempts.
\end{exercise}

\growswhen{A fourth pillar chapter, if FHE-gate / ZK-depth work in the
mid-term trajectory~\cite{helut-trajectory} graduates. Until then it is
Chapter~\ref{ch:open}, not a \texttt{\textbackslash part}.}
